\documentclass[11pt]{article}
\usepackage[review]{acl}
\usepackage{times}
\usepackage{latexsym}
\usepackage[T1]{fontenc}
\usepackage[utf8]{inputenc}
\usepackage{microtype}
\usepackage{inconsolata}
\usepackage{amsmath,amssymb}
\usepackage{booktabs}
\usepackage{array}
\newcommand{\LNGF}{LINGUAFORCE}
\title{\LNGF{}: A Human-Annotated Benchmark\\for Pressure and Influence Strategies in Dialogue}
\author{Anonymous ACL submission}
\begin{document}
\maketitle

\begin{abstract}
Pressure in dialogue can take the form of a threat, an appeal to duty, or an
attempt to make someone feel guilty. A binary label identifies its presence
but does not describe how it is expressed. \LNGF{} provides human annotations
for 4,066 English dialogues at four levels: pressure presence, overall
intensity, fifteen non-exclusive influence strategies, and seven dimensions
of expressed pressure. Ordinary requests and reasoned persuasion are included
to distinguish influence from negative pressure. Twenty annotators completed
the annotation batches. The dataset contains 3,144 dialogues labeled as
exerting pressure and 922 labeled as not exerting pressure. Its labels
support four prediction tasks. On 634 held-out dialogues, three zero-shot
language models and two supervised lexical baselines obtain pressure AUROC
between 0.947 and 0.958. Strategy macro-F1 remains between 0.570 and 0.610,
and ordinal prediction separates the methods more clearly than detection.
An additional matched Qwen3-4B fine-tuning diagnostic produces only small,
task-specific changes, highlighting the difficulty of the fine-grained labels.
The benchmark supports evaluation of both pressure recognition and its
fine-grained description.
\end{abstract}

\section{Introduction}
\emph{Could you stay a little longer?} and \emph{If you leave, you clearly do
not care about this family} ask for the same action. The second makes refusal
carry a moral and emotional cost. Identifying this difference requires
attention to the requested action, the surrounding exchange, and the
consequences attached to refusal.

This distinction also concerns language as social action: an utterance can
shape what a listener is expected or permitted to do, rather than merely
describe a state of affairs \citep{ahearn2001}.

Dialogue datasets study related phenomena such as mental manipulation,
condescension, and toxicity \citep{mentalmanip,multimanip,talkdown2023,toxicchat2023}.
Toxic language is not necessary for pressure, and a direct request is not
necessarily coercive. A benchmark of conversational pressure therefore
needs labels that describe both whether pressure is present and how it is
expressed.

\LNGF{} pairs four kinds of human judgments. A binary label records whether
a dialogue exerts negative pressure. An ordinal label records its overall
intensity. Strategy labels describe the means used to influence the listener.
Seven dimension scores describe directive force, option constraint, normative
pressure, emotional pressure, deceptiveness, toxicity, and explicitness.
A strategy identifies a kind of behavior, while a dimension records its
degree or manner of expression.

The scope is pressure expressed in text. The annotations do not establish
whether a listener actually complied or experienced harm. Ordinary requests,
advice, and reasoned persuasion are included as controls; beneficial influence
such as encouragement is not a separate target of this release.

The contributions are a human-annotated resource with four related label
sets, operational definitions that distinguish pressure from ordinary
influence, and an evaluation design covering detection, intensity, strategy
recognition, and dimension prediction. Evaluation compares three zero-shot
language models with two supervised lexical baselines and includes a matched
before/after fine-tuning diagnostic. Strong detection scores coexist with
lower strategy recognition scores, motivating evaluation beyond the presence
of pressure alone.

\section{Related Work}
\paragraph{Influence strategies.}
Compliance-gaining research studies how speakers seek cooperation
\citep{marwell1967,wiseman1981}. Reciprocity, authority, and social proof
describe common sources of influence \citep{cialdini1984}. Speech-act theory
distinguishes requests, orders, and promises \citep{austin1962,searle1969}.
Normative and moral appeals have also been studied as moral-foundation
signals in language and text classification \citep{graham2009,johnson2018}.
These accounts motivate the strategy inventory. They do not imply that every
request, promise, or appeal to authority is harmful.

\paragraph{Language datasets.}
MentalManip and MultiManip study manipulation in conversations
\citep{mentalmanip,multimanip}. Model-assisted augmentation and distillation
have also been explored for mental-manipulation detection
\citep{reament}. SemEval-2023 Task 3 annotates persuasion techniques in
online news \citep{piskorski2023}. TalkDown and ToxicChat address
condescension and toxicity \citep{talkdown2023,toxicchat2023}, while work on
microaggressions illustrates the difficulty of locating subtle interpersonal
harm in text \citep{breitfeller2019}.
\LNGF{} focuses on dialogue-level pressure and pairs strategy categories
with ordinal judgments. Its evaluation asks whether models can recover
these human judgments from dialogue text.

\paragraph{Modeling choices.}
Contextual encoders such as BERT and RoBERTa are standard supervised text
baselines \citep{devlin2019,liu2019}. We use TF--IDF models instead to make
the lexical-signal comparison explicit and to keep the supervised baseline
separate from the prompted language-model evaluation.

\section{Annotation Scheme}
\subsection{Pressure and Intensity}
Negative pressure is an attempt to obtain an action or decision by making
refusal costly, unacceptable, or difficult to consider. Evidence includes
threatened consequences, guilt, obligations used as leverage, misleading
claims, and hostility directed toward obtaining compliance. Information
exchange, ordinary advice, and disagreement do not automatically qualify.
Annotators consider the complete dialogue rather than individual keywords.
A listener's resistance does not make the pressure disappear.

The binary label records presence or absence. Overall intensity uses six
levels: none (0), very slight (1), weak (2), moderate (3), strong (4), and
very strong (5). Intensity concerns the degree of pressure, not the number
of strategies. In the completed annotations, absence corresponds to
intensity 0 and presence to intensity 1--5.

\subsection{Influence Strategies}
Table~\ref{tab:strategies} lists fifteen non-exclusive labels. Multiple
strategies may occur in one dialogue. Requests and rational persuasion are
included because the task must distinguish ordinary influence from pressure.
A strategy label alone does not determine the binary label.

\begin{table*}[t]
\centering\small
\begin{tabular}{@{}lp{8.5cm}r@{}}
\toprule
\textbf{Strategy} & \textbf{Annotation focus} & \textbf{Count}\\
\midrule
Request/Advice & An open request, suggestion, or recommendation. & 1,644\\
Directive/Order & An instruction that expects compliance. & 2,078\\
Rational Persuasion & Reasons or evidence offered to support an action. & 506\\
Promise/Inducement & A benefit, reward, or favor offered in return for action. & 258\\
Threat/Warning & An adverse consequence used to obtain action. & 705\\
Moral Appeal/Guilt & An appeal to conscience or guilt. & 1,851\\
Obligation/Duty & A role or responsibility invoked to obtain compliance. & 1,712\\
Value Judgement/Shaming & A negative evaluation of the listener's character. & 1,336\\
Authority & Rank, expertise, or status used to demand compliance. & 340\\
Conformity/Social Proof & Others' behavior or agreement used as leverage. & 893\\
Reciprocity/Debt & A past favor invoked to demand repayment through action. & 737\\
Emotional Manipulation & Emotional vulnerability exploited to change a decision. & 882\\
Deception/Gaslighting & Facts or memory distorted to influence the listener. & 30\\
Logical Trap/False Dilemma & A misleading argument or false restriction of alternatives. & 234\\
Verbal Abuse/Toxicity & Insults or hostility used to obtain compliance. & 947\\
\bottomrule
\end{tabular}
\caption{Strategy definitions and positive counts across 4,066 dialogues.
Labels may overlap, so counts do not sum to the number of dialogues.}
\label{tab:strategies}
\end{table*}

Gaslighting is included under deception.\footnote{Gaslighting involves
denial or distortion that leads someone to doubt their memory or judgment.
A disagreement about memory alone is insufficient evidence.}
Deception requires contextual evidence of misleading facts, claims, or
promises. An unverified assertion or an ordinary future promise is not
automatically deceptive.

\subsection{Seven Dimensions}
Each dimension receives an ordinal score of 0 (absent), 1 (low), 2 (moderate),
or 3 (high). Table~\ref{tab:dimensions} defines the seven judgments.

\begin{table}[t]
\centering\small
\begin{tabular}{@{}p{2.6cm}p{4.65cm}@{}}
\toprule
\textbf{Dimension} & \textbf{Judgment}\\
\midrule
Directive force & Strength of the push toward an action.\\
Option constraint & How much the listener's alternatives are narrowed.\\
Normative pressure & Leverage through duty, morality, or social norms.\\
Emotional pressure & Guilt, fear, shame, or pity used as leverage.\\
Deceptiveness & Misleading facts or promises supported by context.\\
Toxicity & Insults, hostility, or verbal abuse.\\
Explicitness & How directly the action or pressure is expressed.\\
\bottomrule
\end{tabular}
\caption{Human-rated dimensions, each scored from 0 to 3.}
\label{tab:dimensions}
\end{table}

These scores are not interchangeable with overall intensity. A request can
be explicit while leaving refusal available. An ultimatum can restrict
options without containing insults. Normative and emotional pressure can
co-occur in the same appeal. The scheme permits these overlaps and does not
assume statistical independence between dimensions.

\section{Dataset and Human Annotation}
\subsection{Composition}
The dataset contains 4,066 English dialogues. The current partition has
3,432 development-pool dialogues and 634 test dialogues
(Table~\ref{tab:composition}). A fixed, intensity-stratified split divides
the development pool into 3,088 training and 344 validation dialogues.

\begin{table}[t]
\centering\small
\setlength{\tabcolsep}{3pt}
\begin{tabular}{@{}lrrr@{}}
\toprule
\textbf{Partition} & \textbf{Total} & \textbf{Pressure} & \textbf{No pressure}\\
\midrule
Development pool & 3,432 & 2,640 & 792\\
Test & 634 & 504 & 130\\
\midrule
All & 4,066 & 3,144 & 922\\
\bottomrule
\end{tabular}
\caption{Counts from the completed human annotation sheets. The development
pool is reserved for training and validation.}
\label{tab:composition}
\end{table}

The dialogues were manually collected, organized, and written in English
around common everyday interaction scenarios. They are not a corpus of
verbatim real-world chat logs. This scenario-based collection provides
examples of different influence strategies, rather than a representative
sample of their frequency in everyday conversation. Dialogue construction
and the annotation procedure below are separate stages.

\subsection{Annotation Procedure}
Twenty computer science graduate students, all proficient in reading
English dialogue, volunteered as annotators. All read the same annotation
guidelines and completed their assigned batches independently without
consulting other annotators' labels. Each dialogue received one annotation
across twenty non-overlapping batches. Each sheet
contained the dialogue and fields for the binary label, intensity, seven
dimensions, fifteen strategies, and optional notes. The annotation package
did not include model predictions. Its instructions asked annotators to read
the full dialogue, assign presence and intensity, rate the dimensions, and
select all applicable strategies.

\subsection{Quality Checks}
The completed sheets were checked for coverage, valid ranges, repeated
records, and consistency between binary presence and intensity. All 4,066
records matched uniquely to their source dialogues using identifier and
text. All required cells are populated and in range. No binary--intensity
contradictions were found. No exact duplicate dialogues were found after
lowercasing and whitespace normalization. These checks detect formatting
and consistency problems; they do not measure agreement between people.
No dialogues have yet received independent duplicate annotations, so
inter-annotator agreement cannot currently be estimated.

\section{Dataset Characteristics}
Pressure is present in 3,144 dialogues (77.3\%). Table~\ref{tab:intensity}
shows the intensity distribution. Moderate pressure is the most frequent
level. Majority-label baselines and metrics sensitive to class imbalance
are therefore necessary for interpreting predictive performance.

\begin{table}[t]
\centering\small
\begin{tabular}{@{}lrrrrrr@{}}
\toprule
Intensity & 0 & 1 & 2 & 3 & 4 & 5\\
\midrule
Count & 922 & 232 & 511 & 1,460 & 775 & 166\\
\bottomrule
\end{tabular}
\caption{Distribution of human-rated overall intensity.}
\label{tab:intensity}
\end{table}

Strategy frequency varies substantially (Table~\ref{tab:strategies}).
Directives and moral appeals are common, whereas contextual evidence for
deception is rare. Micro-averaged performance alone would obscure errors
on rare labels, so evaluation includes macro-F1 and per-label support.
Dimension evaluation likewise reports each dimension rather than only an
average. These frequencies describe this collection, not the prevalence
of such behaviors in everyday conversation.

\section{Benchmark Tasks and Evaluation}
\subsection{Prediction Tasks}
Every task takes dialogue text as input. Human labels are evaluation
targets and are not supplied as test-time features.
\paragraph{Pressure detection.}
Predict binary pressure. Report positive-class F1 and AUROC from a model's
pressure score. Binary decisions are evaluated directly.
\paragraph{Intensity prediction.}
Predict a 0--5 level. Report mean absolute error (MAE) and quadratic-weighted
$\kappa$ (QWK). MAE measures error size, while QWK penalizes larger ordinal
disagreements.
\paragraph{Strategy recognition.}
Predict fifteen non-exclusive labels. Report macro- and micro-F1, with
per-strategy precision, recall, F1, and support in supplementary results.
\paragraph{Dimension prediction.}
Predict seven 0--3 ratings. Report MAE and QWK per dimension and their
unweighted averages. Undefined QWK in degenerate evaluation subsets must
be reported rather than replaced with an arbitrary value.

\subsection{Baselines}
The comparison includes GPT-6 Astra, DeepSeek-Flash, and
Qwen3.8-Max, accessed through a company gateway, the DeepSeek API, and
DashScope, respectively. Each receives the same label
definitions, dialogue text, and output schema without labeled examples.
The schema requests a binary decision and pressure score, intensity,
strategy labels, and dimension ratings. Outputs are checked against
the allowed ranges, with invalid outputs and retries reported.

Supervised baselines use word unigram/bigram or character 3--5-gram TF--IDF
features, capped at 50,000 features with minimum document frequency 2 and
sublinear term frequency. Logistic regression predicts pressure and fifteen
independent strategy labels; ridge regression predicts intensity and each
dimension. Ordinal regression outputs are clipped to the allowed range and
rounded. The constant baseline ignores dialogue text, predicting training
majorities for binary labels and training medians for ordinal labels.
All fitted components use training data only.

Qwen and DeepSeek use temperature 0 with thinking disabled. GPT uses the
gateway's default sampling settings and a fixed low reasoning budget. Because
chain-of-thought prompting can change generated answers \citep{wei2022},
the reported comparison fixes the output instruction and does not treat
hidden reasoning as an evaluation target. These three API models receive a
1,600-token output budget; GPT's budget also includes reasoning tokens.
Technical reports describe the broader OpenAI, Qwen, and DeepSeek model
families \citep{openai2023,qwen2025,deepseek2024}; the names here are provider
aliases, and those reports do not identify the exact deployed versions.
Returned identifiers and usage are saved.
These provider-specific configurations are not matched reasoning budgets.

\paragraph{Matched fine-tuning diagnostic.}
To test whether the annotations support supervised adaptation, we fine-tuned
Qwen3-4B on the 3,088 training dialogues with LoRA. The run used rank 16,
scaling factor 32, dropout 0.05, learning rate $10^{-4}$, three epochs,
batch size 1 with gradient accumulation 8, and a maximum sequence length of
2,048 tokens. The resulting checkpoint and the original Qwen3-4B received
the same 634 test dialogues, output schema, and \texttt{/no\_think} suffix.
This matched comparison is reported as a diagnostic; it is separate from the
zero-shot and lexical baseline comparison.

\subsection{Experimental Controls}
The split and lexical baselines use seed 2026. Validation AUROC selects
pressure regularization from $C\in\{0.1,1,10\}$; validation F1 selects its
decision threshold from 0.10 to 0.90 in increments of 0.01. Strategy
macro-F1 selects $C$ from the same grid and a shared threshold from
$\{0.3,0.4,0.5,0.6\}$. Each ordinal target selects ridge regularization
from $\{0.1,1,10\}$ by validation MAE.

The three API baselines share an English user prompt and test dialogues.
Invalid outputs are retained, with validity evaluated separately
for each task. Out-of-range language-model ratings are not clipped or
silently replaced. Metrics on valid outputs must be read together with
their coverage; comparisons on differing subsets do not establish a
method ranking. Syntax-only normalization closes one missing outer JSON
brace in nine GPT answers and maps one unambiguous strategy-key spelling
in 146 Qwen answers. These rules apply to all three API baselines and leave label values
unchanged; raw responses remain available. Transient GPT transport errors are retried up to three
times with increasing delays; invalid generated answers are not retried.
Reproduction materials include split identifiers, prompts, settings,
predictions, and scoring code. Lexical baseline results use one fixed seed.
We estimate 95\% confidence intervals using 1,000 paired dialogue-level
bootstrap resamples (seed 2026). These quantify test-sample uncertainty
for fixed predictions, not variation across training or API runs.
Appendix~\ref{sec:uncertainty} reports the full intervals and a diagnostic
of false positives on non-pressuring requests and persuasion.

\section{Results and Analysis}
\subsection{Pressure Detection}
\begin{table}[t]
\centering\small
\setlength{\tabcolsep}{3pt}
\begin{tabular}{@{}lrrr@{}}
\toprule
Model & Valid $n$ & AUROC & F1\\
\midrule
GPT-6 Astra & 634 & 0.947 & 0.959\\
Qwen3.8-Max & 634 & 0.958 & 0.953\\
DeepSeek-Flash & 634 & 0.954 & 0.950\\
Word TF--IDF & 634 & 0.949 & 0.942\\
Character TF--IDF & 634 & 0.948 & 0.943\\
Constant & 634 & 0.500 & 0.886\\
\bottomrule
\end{tabular}
\caption{Pressure detection. GPT, Qwen, and DeepSeek are zero-shot; TF--IDF models use human training labels. Constant predicts the training-majority label. Valid counts are out of 634; AUROC additionally requires a valid pressure score (see released metrics).}
\label{tab:pressure_results}
\end{table}
Table~\ref{tab:pressure_results} reports detection on all 634 test dialogues.
Qwen has the highest numerical AUROC (0.958), while GPT has the highest F1
(0.959). Word and character TF--IDF reach AUROC 0.949 and 0.948. For each
language model, the paired 95\% bootstrap interval for the AUROC difference
from word TF--IDF includes zero. The data therefore support strong lexical
predictability, without establishing a clear ranking advantage for prompted
language models. GPT's F1 difference from word TF--IDF is 0.017, with a
paired interval of [0.005, 0.031] under the reported decision rules.

\paragraph{Requests without pressure.}
Of the 130 no-pressure test dialogues, 126 carry human request/advice or
rational-persuasion labels. Negative examples therefore include ordinary
influence attempts, rather than predominantly casual conversation.
The prompted and trained models misclassify 22.2\%--34.1\% of these
dialogues (Appendix~\ref{sec:uncertainty}, Table~\ref{tab:request_controls}),
showing that distinguishing pressure from ordinary influence remains
challenging.

\subsection{Intensity and Strategy Prediction}
\begin{table}[t]
\centering\small
\setlength{\tabcolsep}{3pt}
\begin{tabular}{@{}lrrr@{}}
\toprule
Model & Valid $n$ & MAE $\downarrow$ & QWK $\uparrow$\\
\midrule
GPT-6 Astra & 634 & 0.610 & 0.791\\
Qwen3.8-Max & 634 & 0.891 & 0.714\\
DeepSeek-Flash & 634 & 0.997 & 0.630\\
Word TF--IDF & 634 & 0.683 & 0.736\\
Character TF--IDF & 634 & 0.642 & 0.765\\
Constant & 634 & 1.114 & 0.000\\
\bottomrule
\end{tabular}
\caption{Overall intensity prediction on valid 0--5 outputs. Counts are out of 634. The constant predictor uses the training median.}
\label{tab:intensity_results}
\end{table}
GPT achieves intensity MAE 0.610 and QWK 0.791, compared with 0.642 and
0.765 for character TF--IDF (Table~\ref{tab:intensity_results}). Their MAE
difference is not conclusive (paired 95\% CI: $[-0.104, 0.036]$).
Qwen and DeepSeek have larger intensity errors despite similar pressure
detection scores, indicating that detecting pressure and estimating its
degree pose different challenges.

\begin{table}[t]
\centering\small
\setlength{\tabcolsep}{3pt}
\begin{tabular}{@{}lrrr@{}}
\toprule
Model & Valid $n$ & Macro-F1 & Micro-F1\\
\midrule
GPT-6 Astra & 634 & 0.610 & 0.685\\
Qwen3.8-Max & 634 & 0.590 & 0.667\\
DeepSeek-Flash & 634 & 0.595 & 0.677\\
Word TF--IDF & 634 & 0.573 & 0.738\\
Character TF--IDF & 634 & 0.570 & 0.732\\
Constant & 634 & 0.046 & 0.231\\
\bottomrule
\end{tabular}
\caption{Fifteen-label strategy recognition. A task-valid response supplies all fifteen binary labels. One unambiguous key spelling is normalized without changing its value.}
\label{tab:strategy_results}
\end{table}

GPT leads strategy macro-F1 (0.610), while word TF--IDF leads micro-F1
(0.738; Table~\ref{tab:strategy_results}). Both differences between these
models persist under paired resampling (Appendix~\ref{sec:uncertainty}).
The ranking depends on whether strategies receive equal weight or their
decisions are pooled, motivating both metrics for this uneven inventory.
Per-strategy scores and support counts are available in the supplement.

\subsection{Dimension Prediction and Output Validity}
\begin{table}[t]
\centering\small
\setlength{\tabcolsep}{3pt}
\begin{tabular}{@{}lrrr@{}}
\toprule
Model & Valid $n$ & MAE $\downarrow$ & QWK $\uparrow$\\
\midrule
GPT-6 Astra & 634 & 0.430 & 0.606\\
Qwen3.8-Max & 295 & 0.631 & 0.540\\
DeepSeek-Flash & 634 & 0.636 & 0.547\\
Word TF--IDF & 634 & 0.422 & 0.631\\
Character TF--IDF & 634 & 0.424 & 0.614\\
Constant & 634 & 0.760 & 0.000\\
\bottomrule
\end{tabular}
\caption{Dimension prediction on complete, range-valid seven-dimensional outputs. MAE and QWK are averaged over dimensions. Different coverage prevents a direct ranking; Table~\ref{tab:common_dimensions} uses identical dialogues.}
\label{tab:dimension_results}
\end{table}
GPT, DeepSeek, and the supervised baselines return range-valid dimension
profiles for all 634 dialogues after syntax normalization. Qwen does so for
295 dialogues (46.5\%); 339 contain at least one rating outside 0--3.
Those values are retained and counted as invalid, rather than clipped.
Table~\ref{tab:dimension_results} reports valid-profile metrics and coverage.

\begin{table}[t]
\centering\small
\setlength{\tabcolsep}{3pt}
\begin{tabular}{@{}lrrr@{}}
\toprule
Model & MAE $\downarrow$ & QWK $\uparrow$ & All exact $\uparrow$\\
\midrule
GPT-6 Astra & 0.393 & 0.595 & 0.667\\
Qwen3.8-Max & 0.631 & 0.540 & 0.331\\
DeepSeek-Flash & 0.593 & 0.524 & 0.506\\
Word TF--IDF & 0.412 & 0.656 & 0.618\\
Character TF--IDF & 0.408 & 0.582 & 0.615\\
Constant & 0.847 & 0.000 & 0.471\\
\bottomrule
\end{tabular}
\caption{Aligned dimension comparison. MAE/QWK use the same 295 dialogues with valid profiles from every model. All exact instead uses all 634 dialogues and seven ratings each, counting missing or out-of-range ratings as incorrect. The common subset is selected by model output validity and may be easier than the full test set.}
\label{tab:common_dimensions}
\end{table}

Table~\ref{tab:common_dimensions} compares all methods on the same 295
dialogues. GPT has the lowest MAE (0.393), while word TF--IDF has the highest
mean QWK (0.656). This subset is selected by output validity and need not
represent the full test set. To retain the full denominator, the table also
reports exact rating accuracy over all $634\times7$ judgments, counting
invalid ratings as incorrect. GPT obtains 0.667 and word TF--IDF 0.618;
Qwen obtains 0.331 under this fixed prompting configuration.
On the full test set, GPT and word TF--IDF have no clear difference in
dimension MAE (paired 95\% CI: $[-0.022, 0.036]$), while GPT has higher
exact rating accuracy (difference 0.049; CI: $[0.026, 0.071]$).

Across tasks, no method leads every metric. The labels support distinctions
between detecting pressure, naming its strategy, and estimating its degree.
These cross-model comparisons do not isolate the benefit of training a given
language model; the following diagnostic compares Qwen3-4B before and after training.

\subsection{Matched Fine-Tuning Diagnostic}
Table~\ref{tab:finetune_4b} compares the original Qwen3-4B with the
checkpoint fine-tuned on the training split. The two runs use the same test
inputs and decoding instruction. The reported valid $n$ is the common number
of dialogues with task-valid outputs; no missing or out-of-range values are
repaired. ``SFT'' denotes the LoRA fine-tuned checkpoint.

\begin{table*}[t]
\centering\small
\setlength{\tabcolsep}{4pt}
\begin{tabular}{@{}llrrrr@{}}
\toprule
\textbf{Task} & \textbf{Metric} & \textbf{Base} & \textbf{SFT} &
\textbf{SFT--Base} & \textbf{Paired 95\% CI}\\
\midrule
Pressure & F1 & 0.887 & 0.887 & $-0.0002$ & $[-0.0030,\ 0.0024]$\\
Intensity & MAE $\downarrow$ & 1.943 & 1.940 & $-0.0032$ & $[-0.0126,\ 0.0063]$\\
Intensity & QWK $\uparrow$ & 0.214 & 0.216 & $+0.0016$ & $[-0.0036,\ 0.0073]$\\
Strategies & Micro-F1 $\uparrow$ & 0.588 & 0.586 & $-0.0015$ & $[-0.0040,\ 0.0010]$\\
Strategies & Macro-F1 $\uparrow$ & 0.496 & 0.495 & $-0.0011$ & $[-0.0042,\ 0.0024]$\\
Dimensions & MAE $\downarrow$ & 0.878 & 0.874 & $-0.0039$ & $[-0.0096,\ 0.0011]$\\
Dimensions & Mean QWK $\uparrow$ & 0.305 & 0.312 & $+0.0069$ & $[0.0009,\ 0.0135]$\\
\bottomrule
\end{tabular}
\caption{Matched Qwen3-4B fine-tuning comparison on the held-out test set.
The paired valid sample sizes are 634 for pressure and intensity, 623 for
strategies, and 622 for dimensions. Intervals use 1,000 paired dialogue-level
bootstrap resamples; they quantify uncertainty from the observed test sample
for this fixed training run.}
\label{tab:finetune_4b}
\end{table*}

Fine-tuning changes little across the four tasks. Pressure F1 is unchanged,
intensity is nearly unchanged, and both strategy scores are slightly lower.
Mean dimension QWK increases from 0.305 to 0.312 on the common valid subset,
but the absolute gain is small. This result does not support a claim of broad
performance improvement from fine-tuning; it instead shows that the
fine-grained labels remain difficult even when the model is adapted to the
training split.



\section{Conclusion}
\LNGF{} provides human judgments of pressure presence, intensity, strategies,
and seven dimensions for 4,066 dialogues. Its labels support evaluation
beyond binary detection, including whether models identify the kind and
degree of pressure expressed in context. Three zero-shot language models
and supervised lexical baselines perform strongly on binary detection,
while strategy and dimension results reveal different strengths and output
reliability. A matched Qwen3-4B fine-tuning run produces no broad gain across
tasks, with only a small improvement in mean dimension QWK. This result
supports using the resource to measure fine-grained understanding rather than
assuming that adaptation makes every task easy. Independent annotation
agreement remains necessary to assess the reliability of the human reference
labels.

\section*{Limitations}
The benchmark covers English dialogue and negative pressure. It does not
measure listener compliance, psychological harm, or positive influence.
The strategy inventory and scales are operational choices whose reliability
requires human agreement evidence.

Some labels have few positive examples, limiting the precision of per-label
evaluation. The current duplicate check detects exact normalized matches
but does not rule out paraphrases or shared scenario templates.
The fine-tuning comparison uses one Qwen3-4B run and one fixed configuration;
it does not estimate variation across random seeds, model sizes, or training
recipes.
Generalization to other languages, settings, and naturally occurring
interactions requires additional evaluation.
The comparison also does not cover other foundation-model families, such as
Llama or GLM, whose behavior may differ under the same tasks
\citep{touvron2023,glm2024}.

\section*{Ethical Considerations}
The dialogues contain threats, shaming, deception, and verbal abuse.
Predictions should not be used to diagnose individuals or infer intent or
harm beyond the text. Annotators and data users should be informed about
potentially distressing content.
Participation was voluntary and unpaid.
The dataset creators permit release of the dialogue texts and annotations.
The anonymous supplement provides the data and experimental materials for
review. The dialogues concern everyday scenarios and should not be treated
as verified accounts of identifiable individuals.
Language-model assistance was used for manuscript drafting and editing
and for preparing experimental code. The benchmark reference labels were
provided by human annotators.

\bibliography{references}

\appendix
\section{Uncertainty Estimates}
\label{sec:uncertainty}
\paragraph{Paired comparisons.}
For strategy recognition, GPT minus word TF--IDF has a 95\% interval
of $[0.012, 0.062]$ for macro-F1 and $[-0.068, -0.037]$ for micro-F1.
For dimension MAE, the corresponding interval includes zero on both
the full test set and the shared valid-profile subset. These comparisons
support task- and metric-specific conclusions rather than a single
overall model ranking.

\paragraph{Interval estimation.}
All intervals are pointwise 95\% intervals, without a multiple-comparison
adjustment. Bootstrap intervals use the 2.5th and 97.5th percentiles of
1,000 dialogue resamples with replacement. Differences use paired draws,
not subtraction of marginal interval endpoints. They are conditional on
the observed split and predictions and do not establish stability across
training seeds, API runs, or unseen settings.

For the API and lexical baselines, intensity and strategy intervals use all 634 test dialogues
(Table~\ref{tab:ordinal_strategy_ci}). Dimension MAE and mean QWK are
reported both on the shared 295-dialogue valid-profile subset
(Table~\ref{tab:dimension_ci}) and on all 634 dialogues for methods with
complete valid profiles (Table~\ref{tab:dimension_full_ci}). Exact rating
accuracy always uses all 634 dialogues and counts invalid ratings as
incorrect. Resampling whole dialogues preserves dependence among their
seven ratings and fifteen strategy labels.

Mean dimension QWK is undefined if any of its seven component scores is
undefined. Such draws are excluded from that interval and counted in the
tables. Paired QWK differences retain only draws defined for both methods.
The shared subset has sparse dimension categories, so these conditional
QWK intervals should not be interpreted as a stable model ranking.
No undefined value is replaced with zero. All MAE, strategy-F1, and exact
accuracy draws are defined. The supplementary JSON includes paired
differences against both lexical baselines and counts of defined draws.

\paragraph{Non-pressuring influence.}
Among the 130 human-labeled no-pressure test dialogues, 125 carry a
request/advice label and 45 a rational-persuasion label. Their union
contains 126 dialogues, selected independently of model predictions.
Table~\ref{tab:request_controls} reports false-positive rates with Wilson
95\% intervals. The strategy labels are not independent annotations of
action requests, so this diagnostic does not rule out lexical or scenario
shortcuts. Supplementary results provide each group, all 130 negatives,
and false-positive identifiers. No model is refitted for this analysis.

\begin{table}[t]
\centering\small
\setlength{\tabcolsep}{4pt}
\begin{tabular}{@{}lrrr@{}}
\toprule
Model & FP & FPR (\%) & 95\% CI\\
\midrule
GPT-6 Astra & 28 & 22.2 & [15.8, 30.2]\\
Qwen3.8-Max & 35 & 27.8 & [20.7, 36.2]\\
DeepSeek-Flash & 43 & 34.1 & [26.4, 42.8]\\
Word TF--IDF & 38 & 30.2 & [22.8, 38.7]\\
Character TF--IDF & 33 & 26.2 & [19.3, 34.5]\\
Constant & 126 & 100.0 & [97.0, 100.0]\\
\bottomrule
\end{tabular}
\caption{False positives on 126 test dialogues labeled no pressure and request/advice or rational persuasion. Each dialogue is counted once; all systems return valid binary decisions. Intervals are Wilson intervals.}
\label{tab:request_controls}
\end{table}

\begin{table*}[t]
\centering\small
\setlength{\tabcolsep}{4pt}
\begin{tabular}{@{}lrrrr@{}}
\toprule
Model & Intensity MAE & Intensity QWK & Strategy macro-F1 & Strategy micro-F1\\
\midrule
GPT-6 Astra & [0.550, 0.672] & [0.752, 0.827] & [0.592, 0.626] & [0.672, 0.699]\\
Qwen3.8-Max & [0.823, 0.956] & [0.672, 0.753] & [0.572, 0.606] & [0.652, 0.681]\\
DeepSeek-Flash & [0.923, 1.068] & [0.580, 0.678] & [0.578, 0.612] & [0.662, 0.692]\\
Word TF--IDF & [0.631, 0.740] & [0.702, 0.768] & [0.550, 0.594] & [0.722, 0.751]\\
Character TF--IDF & [0.593, 0.696] & [0.732, 0.794] & [0.546, 0.591] & [0.716, 0.747]\\
Constant & [1.024, 1.200] & [0.000, 0.000] & [0.044, 0.049] & [0.217, 0.245]\\
\bottomrule
\end{tabular}
\caption{Pointwise 95\% bootstrap intervals on all 634 test dialogues, with 1,000 paired dialogue resamples. All draws are defined. Point estimates appear in Tables~\ref{tab:intensity_results} and~\ref{tab:strategy_results}.}
\label{tab:ordinal_strategy_ci}
\end{table*}
\begin{table*}[t]
\centering\small
\setlength{\tabcolsep}{4pt}
\begin{tabular}{@{}lrrrr@{}}
\toprule
Model & MAE (295) & QWK (295) & Defined QWK draws & Exact accuracy (634)\\
\midrule
GPT-6 Astra & [0.352, 0.437] & [0.546, 0.635] & 942 & [0.651, 0.683]\\
Qwen3.8-Max & [0.578, 0.686] & [0.487, 0.586] & 1000 & [0.311, 0.354]\\
DeepSeek-Flash & [0.539, 0.656] & [0.465, 0.572] & 1000 & [0.486, 0.527]\\
Word TF--IDF & [0.377, 0.446] & [0.548, 0.691] & 876 & [0.600, 0.635]\\
Character TF--IDF & [0.376, 0.440] & [0.542, 0.615] & 628 & [0.596, 0.632]\\
Constant & [0.798, 0.897] & [0.000, 0.000] & 628 & [0.453, 0.489]\\
\bottomrule
\end{tabular}
\caption{Dimension 95\% intervals. MAE/QWK use the same 295 complete-profile dialogues; exact rating accuracy uses all 634 dialogues, counting invalid outputs as wrong. QWK intervals exclude draws in which any dimension has undefined QWK; the resulting conditional intervals require caution. MAE and accuracy use all 1,000 draws.}
\label{tab:dimension_ci}
\end{table*}
\begin{table*}[t]
\centering\small
\setlength{\tabcolsep}{4pt}
\begin{tabular}{@{}lrrr@{}}
\toprule
Model & MAE & QWK & Defined QWK draws\\
\midrule
GPT-6 Astra & [0.402, 0.459] & [0.571, 0.640] & 1000\\
DeepSeek-Flash & [0.600, 0.671] & [0.515, 0.578] & 1000\\
Word TF--IDF & [0.399, 0.445] & [0.584, 0.675] & 1000\\
Character TF--IDF & [0.403, 0.445] & [0.580, 0.661] & 999\\
Constant & [0.726, 0.789] & [0.000, 0.000] & 993\\
\bottomrule
\end{tabular}
\caption{Dimension 95\% intervals for methods with complete valid profiles on all 634 dialogues. Qwen is excluded from this full-profile comparison. Undefined QWK draws are excluded and counted; MAE uses all 1,000 draws.}
\label{tab:dimension_full_ci}
\end{table*}
\end{document}
